\begin{otherlanguage}{english}
\chapter*{\centering Abstract}
Retail networks must choose replenishment policies for heterogeneous demand. This dissertation proposes the Adaptive Inventory Policy Engine (AIPE), a methodology comprising five stages: specifying the decision, representing context, building comparable operational evidence, recommending a policy, and evaluating its consequences. Its Policy Selection Engine (PSE) supports two approaches: mapping a supplied profile to a policy and estimating a recommendation for each series. The methodology distinguishes policy-family selection from parameter calibration and allows a reference policy to be retained when evidence for replacing it is insufficient. Seven policies are evaluated in two separate applications. In Bahia, 145 series at two origins yield 4,350 fits; profile-based selection increases mean cost by 1.93\% and mean fill rate by 3.89 percentage points relative to global selection. In the M5 pilot, 30 series from ten stores at three origins yield 1,350 fits. The original predictive selector has higher cost and lower service than the reference. An exploratory comparison reusing these simulations evaluates selection based on predicted cost and service: with fallback, mean cost decreases by 1.56\% and service increases by 0.27 percentage points, but only three decisions change and one series accounts for approximately 92\% of savings. Generalizable superiority is not established. Retrospective oracles distinguish selection errors, portfolio gaps, and physical infeasibility. The contribution is the formulation and implementation of the AIPE procedure, supported by evidence about its utility and limitations. Retrospective cohorts, different calibrations, short horizons, and development on the pilot require independent confirmation.

\vspace{2em}
\textbf{Keywords}: Selection methodology; Inventory management; Contextual policy selection; Service level; AIPE.
\end{otherlanguage}
